\documentclass[12pt]{article}
\usepackage[utf8]{inputenc}
\usepackage[T1]{fontenc}
\usepackage{hyperref}
\title{Mokati Language Resources for K-12 Curriculum Development}
\author{Compiled by Claude Sonnet 5.5 (claude-sonnet-5-5) for Andr\'as Kornai}
\date{}
\begin{document}
\maketitle

\section{Introduction}
Mokati (ISO 639-3 wnb), also called Wanambre, is spoken in Sumgilbar Rural LLG, Madang Province, Papua New Guinea. Wikipedia reports about 590 speakers (2003); Global Recordings Network gives 490. There are two main dialect areas, centred on Wanambre and Mawet, and on Tinami and Kotakot. Mokati should not be confused with Wanap (wnp), a Torricelli language of Sandaun Province. No school use was found.

\section{Primary Educational Resources}
None found.

\section{Linguistic and Reference Materials}
\begin{itemize}
\item Wikipedia, Mokati language: location, dialects and sample vocabulary. \url{https://en.wikipedia.org/wiki/Mokati_language}
\end{itemize}
No dictionary or grammar was found.

\section{Religious and Community Resources}
None found.

\section{Digital and Audio Resources}
Global Recordings Network has a language page for wnb. The search did not show what recordings, if any, it holds. \url{https://globalrecordings.net/en/language/wnb}

\section{Curriculum Development Context}
Almost nothing is published on Mokati. Any curriculum would first require documentation work with speakers.

\end{document}